% Transposition surgery, Euler's inequality, and two-sided faces.
%
% Self-contained map-theoretic foundations for the compositional route.
% Every result carries an inline verification-status marker; the appendix
% ledger is generated and checked by tools/check_status.py.
%
% This file is \input by the companion and is maintained separately from it.

\section{Foundations: transposition surgery and two-sided faces}
\label{sec:map-euler-foundations}

The compositional route needs one classical fact at its classical boundary:
in a bridgeless spherical cubic map, the two sides of every edge lie on
\emph{different} faces. In the informal literature this is treated as a
consequence of the Jordan curve theorem and therefore as visually obvious.
It is obvious. It is also, in a permutation model of maps, a corollary of
Euler's inequality, and that route is entirely finite and combinatorial.

This section gives the finite argument in full. The mathematics is
standard --- the surgery lemma is folklore and Euler's inequality for
hypermaps goes back to Jacques --- but the path is unusually short when the
genus inequality is proved \emph{first}, and it needs remarkably little: no
disjointness of the transpositions, no orientability argument, and no
Jordan curve theorem.

\subsection{Setting}

Throughout, $D$ is a finite set of \emph{darts} and $\Sym(D)$ its symmetric
group. For $\pi \in \Sym(D)$ write $c(\pi)$ for the number of orbits of
$\pi$ on $D$, \emph{counting fixed points as orbits}. For a family
$G \subseteq \Sym(D)$ write $c\langle G\rangle$ for the number of orbits of
the subgroup it generates.

\begin{definition}[combinatorial map]\label{def:map}
A \emph{map} is a triple $(D,\sigma,\alpha)$ with $\sigma \in \Sym(D)$ and
$\alpha \in \Sym(D)$ a fixed-point-free involution. We call $\sigma$ the
\emph{vertex rotation}, $\alpha$ the \emph{edge flip}, and
$\varphi := \sigma\alpha$ the \emph{face permutation}, whose orbits are the
\emph{faces}. Set
\[
  V := c(\sigma), \qquad E := c(\alpha) = |D|/2, \qquad
  F := c(\varphi), \qquad C := c\langle \sigma,\alpha\rangle .
\]
The map is \emph{connected} if $C = 1$, and \emph{spherical} if it is
connected and $V - E + F = 2$.
\end{definition}

Two conventions deserve comment. Faces are orbits of $\sigma\alpha$ rather
than $\alpha\sigma$; the two are conjugate, so no count below is affected.
And an \emph{edge} is a pair $\{d,\alpha d\}$ whose two elements are its
two \emph{sides} --- whether they lie on one face or two is exactly the
question at issue.

\begin{definition}[bridge]\label{def:bridge}
An edge $e = \{d, \alpha d\}$ is a \emph{bridge} if deleting it increases
the number of connected components. Formally, put
$\alpha_0 := \alpha \circ (d\;\alpha d)$, an involution fixing $d$ and
$\alpha d$ and agreeing with $\alpha$ elsewhere, and let
$C_0 := c\langle \sigma, \alpha_0\rangle$. Then $e$ is a bridge exactly
when $C_0 > C$.
\end{definition}

Note that $\alpha_0$ is \emph{not} the edge flip of a map: it has fixed
points. This is intentional and harmless, since none of the counting below
assumes fixed-point-freeness. Retaining the darts $d, \alpha d$ rather than
removing them is what leaves $\sigma$, and hence $V$, literally unchanged
by the deletion, which is what keeps the bookkeeping honest.

\subsection{Transposition surgery}

Everything rests on one lemma about what a transposition does to an orbit
count.

\begin{lemma}[split/merge]\label{lem:split-merge}
\Sverified{Mettapedia.GraphTheory.FourColor.GoertzelV24PermutationOrbitSurgery.orbitCount_swap_mul_of_sameCycle}
Let $\pi \in \Sym(D)$, let $a \neq b$ in $D$, and let $\tau = (a\;b)$. If
$a$ and $b$ lie on the same $\pi$-orbit then $c(\tau\pi) = c(\pi) + 1$;
otherwise $c(\tau\pi) = c(\pi) - 1$.
\end{lemma}

\begin{proof}
Since $(\tau\pi)(x) = \tau(\pi x)$, the permutations $\tau\pi$ and $\pi$
agree at every $x$ with $\pi x \notin \{a,b\}$. Only two points are
affected: the $\pi$-preimages of $a$ and of $b$.

Suppose first that $a,b$ share a $\pi$-orbit, written
\[
  (a,\ \pi a,\ \pi^2 a,\ \dots,\ \pi^{m-1}a), \qquad \pi^m a = a, \qquad
  b = \pi^k a, \quad 0 < k < m .
\]
Then $\pi$ sends $\pi^{k-1}a \mapsto b$ and $\pi^{m-1}a \mapsto a$, so
$\tau\pi$ sends $\pi^{k-1}a \mapsto a$ and $\pi^{m-1}a \mapsto b$. The
single orbit therefore breaks into
$(a, \pi a, \dots, \pi^{k-1}a)$ and $(b, \pi^{k+1}a, \dots, \pi^{m-1}a)$:
two orbits where there was one, all others untouched.

For the second case note $\tau(\tau\pi) = \pi$, and that $a$ and $b$ do
share a $\tau\pi$-orbit --- indeed $\tau\pi$ carries the $\pi$-preimage of
$b$ to $a$ and the $\pi$-preimage of $a$ to $b$, welding the two orbits
into one. The first case applied to $\tau\pi$ then gives
$c(\pi) = c(\tau\pi) + 1$.
\end{proof}

The second case repays attention, because it is the only place a reader
might reach for a picture. The identity $\tau(\tau\pi) = \pi$ says the two
cases are one computation read in opposite directions: a transposition is
its own inverse, so splitting and merging are a single operation seen from
two sides.

\begin{corollary}[right-handed form]\label{cor:split-merge-right}
\Sverified{Mettapedia.GraphTheory.FourColor.GoertzelV24PermutationOrbitSurgery.orbitCount_swap_mul_of_not_sameCycle}
With $\pi, a, b, \tau$ as above, $c(\pi\tau)$ obeys the same dichotomy.
\end{corollary}

\begin{proof}
$\tau(\pi\tau)\tau^{-1} = \tau\pi$, so $\pi\tau$ and $\tau\pi$ are
conjugate and have equal orbit counts; apply
Lemma~\ref{lem:split-merge}.
\end{proof}

We use the right-handed form below, since deleting an edge multiplies the
face permutation on the right.

\begin{remark}
Lemma~\ref{lem:split-merge} is the exact orbit-counting refinement of the
familiar parity statement $\operatorname{sgn}(\tau\pi) =
-\operatorname{sgn}(\pi)$. Parity records that the count changes by an odd
number; the lemma records that it changes by exactly one, and says which
way.
\end{remark}

\subsection{Euler's inequality}

\begin{theorem}[Euler bound]\label{thm:euler-bound}
\Sverified{Mettapedia.GraphTheory.FourColor.GoertzelV24MapEulerBound.orbitCount_add_orbitCount_mul_swapProduct_le}
Let $\sigma \in \Sym(D)$ and let $\tau_1,\dots,\tau_n$ be transpositions,
\emph{not assumed disjoint}, with product $\beta = \tau_1\tau_2\cdots\tau_n$.
Then
\[
  c(\sigma) + c(\sigma\beta) \;\le\; n + 2\,c\langle \sigma, \tau_1, \dots, \tau_n\rangle .
\]
\end{theorem}

\begin{proof}
Induction on $n$.

For $n = 0$ we have $\beta = 1$, the generated group is $\langle\sigma\rangle$
whose orbits are the $\sigma$-orbits, and both sides equal $2c(\sigma)$.
This base case is an \emph{equality}; all the slack in the theorem is
created by the inductive step.

For the step write $\beta = \tau\beta'$ with $\tau = (a\;b)$ and
$\beta' = \tau_2\cdots\tau_n$. Put
$G' := \langle \sigma, \tau_2, \dots, \tau_n\rangle$ and
$G := \langle \sigma, \tau, \tau_2, \dots, \tau_n\rangle$, with
$C' = c\langle G'\rangle$ and $C = c\langle G\rangle$. The inductive
hypothesis reads $c(\sigma) + c(\sigma\beta') \le (n-1) + 2C'$.

The key algebraic step is that inserting $\tau$ on the right of $\sigma$ is
the same as inserting a conjugate transposition on the left of
$\sigma\beta'$:
\[
  \sigma\beta = \sigma\tau\beta'
  = (\sigma\tau\sigma^{-1})(\sigma\beta')
  = \tau'(\sigma\beta'), \qquad \tau' = (\sigma a\;\;\sigma b) .
\]
Since $a \neq b$ and $\sigma$ is injective, $\sigma a \neq \sigma b$, so
Lemma~\ref{lem:split-merge} applies to $\tau'$ and $\sigma\beta'$.

\emph{Case (i): $a$ and $b$ lie in the same $G'$-orbit.} Adjoining $\tau$
joins nothing new, so $C = C'$. The lemma gives
$c(\sigma\beta) \le c(\sigma\beta') + 1$ in either branch, whence
\[
  c(\sigma) + c(\sigma\beta) \le c(\sigma) + c(\sigma\beta') + 1
  \le (n-1) + 2C' + 1 = n + 2C .
\]

\emph{Case (ii): $a$ and $b$ lie in different $G'$-orbits.} Adjoining
$\tau$ fuses exactly those two, so $C = C' - 1$. We claim the lemma must
take its merging branch. Indeed $\sigma \in G'$, so $\sigma a$ and
$\sigma b$ again lie in different $G'$-orbits; and $\sigma\beta' \in G'$,
so $\langle \sigma\beta'\rangle \le G'$ and every $\sigma\beta'$-cycle lies
inside a single $G'$-orbit. Hence $\sigma a$ and $\sigma b$ lie on
different $\sigma\beta'$-cycles and
$c(\sigma\beta) = c(\sigma\beta') - 1$. Therefore
\[
  c(\sigma) + c(\sigma\beta) = c(\sigma) + c(\sigma\beta') - 1
  \le (n-1) + 2C' - 1 = n + 2C . \qedhere
\]
\end{proof}

The two cases repay contrast. In case (i) the component count is frozen
while the face count may rise, so the inequality loses ground --- by
exactly one unit per independent cycle, which is the genus. In case (ii)
the component count falls, but the face count is \emph{forced} to fall with
it and the ledger balances exactly. The bound is tight because a new edge
across two components can never create a face.

\begin{corollary}[Euler's inequality for maps]\label{cor:euler-map}
\Sverified{Mettapedia.GraphTheory.FourColor.GoertzelV24InvolutionEdgeList.orbitCount_add_orbitCount_mul_le'}
For a map $(D,\sigma,\alpha)$ we have $V - E + F \le 2C$, with the edges
taken to be the two-cycles of $\alpha$.
\end{corollary}

\begin{proof}
Present $\alpha$ as an explicit list of pairwise disjoint transpositions by
repeatedly stripping a two-cycle: from a moved point $a$, replace $\alpha$
by $\mathrm{swap}(a, \alpha a)\,\alpha$, which fixes both $a$ and $\alpha a$
and agrees with $\alpha$ elsewhere, so the moved set strictly shrinks and
the recursion terminates. The product of the emitted list is $\alpha$.
Theorem~\ref{thm:euler-bound} applied to that list gives $V + F \le E + 2C$,
since disjointness makes a single $\alpha$-step coincide with a step of one
of the transpositions.
\end{proof}

\begin{remark}[the bookkeeping, discharged]
\Sverified{Mettapedia.GraphTheory.FourColor.GoertzelV24InvolutionEdgeList.edgeList_fst_ne_snd}
That each emitted pair is a genuine two-cycle holds by construction, since a
pair is produced only from a moved point; no involutivity is even needed for
it. With this the map form of the bound is unconditional.
\end{remark}

\begin{remark}\label{rem:genus}
For a connected map this reads $V - E + F \le 2$, and the \emph{genus}
$g := (2 - V + E - F)/2$ is precisely the accumulated slack of case~(i).
Corollary~\ref{cor:euler-map} is the statement $g \ge 0$; spherical maps
are those achieving equality.
\end{remark}

\subsection{Bridges and two-sided faces}

\begin{theorem}\label{thm:same-face-bridge}
\Sverified{Mettapedia.GraphTheory.FourColor.GoertzelV24BridgeTwoSided.two_le_components_of_sameCycle}
Let $(D,\sigma,\alpha)$ be a spherical map and let $e = \{d, \alpha d\}$ be
an edge whose two sides lie on the \emph{same} face. Then $e$ is a bridge.
\end{theorem}

\begin{proof}
Let $\tau := (d\;\;\alpha d)$ and $\alpha_0 := \alpha\tau$ as in
Definition~\ref{def:bridge}, and let $C_0$ be the number of components
after deleting $e$. Three observations.

First, $\sigma$ is untouched, so the vertex count of the deleted structure
is still $V$.

Second, $\alpha_0$ is a product of the $E-1$ disjoint transpositions of
$\alpha$ other than $\tau$, together with two fixed points, and
$\sigma\alpha_0 = \sigma\alpha\tau = \varphi\tau$. Since $d$ and $\alpha d$
lie on a common $\varphi$-orbit, Corollary~\ref{cor:split-merge-right}
gives $c(\sigma\alpha_0) = F + 1$.

Third, Theorem~\ref{thm:euler-bound} applied to $\sigma$ and those $E-1$
transpositions gives $V + c(\sigma\alpha_0) \le (E-1) + 2C_0$, the
generated group having exactly the components of the deleted structure as
in Corollary~\ref{cor:euler-map}.

Combining, $V + F + 1 \le E - 1 + 2C_0$. Sphericity gives $V + F = E + 2$,
so $E + 3 \le E - 1 + 2C_0$, that is $C_0 \ge 2$. The original map was
connected, so deleting $e$ increased the component count and $e$ is a
bridge.
\end{proof}

\begin{corollary}[two-sidedness]\label{cor:two-sided}
\Sverified{Mettapedia.GraphTheory.FourColor.GoertzelV24BridgeTwoSided.not_sameCycle_of_components_lt_two}
In a bridgeless spherical map every edge has two distinct faces on its two
sides.
\end{corollary}

\begin{proof}
Contrapositive of Theorem~\ref{thm:same-face-bridge}.
\end{proof}

\begin{remark}[the converse, and what it actually costs]\label{rem:converse}
\Sprose
The converse --- a bridge has one face on both sides --- is also true, but
it is \emph{not} obtainable from Theorem~\ref{thm:euler-bound}. Running the
same computation with $C_0 = 2$ yields only $c(\varphi\tau) \le F + 1$,
which admits both branches. It was tempting to conclude that the converse
therefore needs the equality case of Euler; it does not. It has an
elementary proof of a different kind, given next, which uses no counting at
all.
\end{remark}

\begin{theorem}[the walk lemma]\label{thm:walk-lemma}
\Sverified{Mettapedia.GraphTheory.FourColor.GoertzelV24BridgeTwoSided.sameCycle_of_bridge}
Let $\varphi$ be a permutation of $D$, let $\mathrm{side} \colon D \to
\{0,1\}$, and let $d, o \in D$. Suppose the first step off $d$ changes
side, and that from any point on the far side other than $o$ the next step
stays on the far side. Then $d$ and $o$ lie on one $\varphi$-orbit.
\end{theorem}

\begin{proof}
The orbit of $d$ is finite and returns to $d$, so among $k \ge 1$ there is a
least one with $\mathrm{side}(\varphi^k d) = \mathrm{side}(d)$. It is not
$k = 1$, since the first step crosses. So $\varphi^{k-1}d$ lies on the far
side by minimality, while $\varphi^{k}d$ does not. The only far-side point
whose next step leaves the far side is $o$; hence $\varphi^{k-1} d = o$.
\end{proof}

Applied to a bridge, with $\mathrm{side}$ recording which component of the
edge-deleted map a dart's vertex lies in, this is precisely the converse: a
facial walk crossing a bridge has exactly one door back, namely the bridge's
other dart. Note what the hypotheses do \emph{not} include --- that $o$
itself lies on the far side is not needed, and the formal statement omits
it.

\begin{theorem}[the converse, packaged]\label{thm:converse-packaged}
\Sverified{Mettapedia.GraphTheory.FourColor.GoertzelV24BridgeTwoSided.sameCycle_of_not_wordReachable}
If deleting an edge separates its two endpoints --- that is, if the edge is
a bridge --- then its two sides lie on one face orbit.
\end{theorem}

\begin{proof}
Take $\mathrm{side}(x)$ to record whether $x$ is reachable from $a$ in the
edge-deleted map. Then $\mathrm{side}(a)$ holds by reflexivity; the first
step off $a$ lands on $\sigma b$, which is reachable from $b$ and hence not
from $a$, so the step crosses; and from any far-side point $x \neq b$ the
edge product cannot carry $x$ onto $a$ or $b$ --- both are fixed by it, and
it is injective --- so the step is a rotation composed with a remaining
flip, neither of which changes the component. Theorem~\ref{thm:walk-lemma}
applies.
\end{proof}

\begin{corollary}[bridges are exactly one-sided edges]\label{cor:bridge-iff}
\Sprose
In a spherical map, an edge is a bridge if and only if its two sides lie on
the same face.
\end{corollary}

\begin{proof}
Theorem~\ref{thm:same-face-bridge} and Theorem~\ref{thm:converse-packaged}.
\end{proof}

Both halves are now finite and combinatorial. The forward direction counts;
the converse walks. Neither looks at a picture.

\subsection{Kempe components and boundary stubs}

One further counting fact belongs here, because the compositional route's
terminal analysis uses it and because it is easy to state in the wrong place.

\begin{lemma}[a Kempe component holds at most two stubs]\label{lem:two-stubs}
\Sverified{Mettapedia.GraphTheory.FourColor.GoertzelV24KempeStubBound.card_stubs_le_two}
Let $H$ be a finite graph in which every vertex has degree at most two, and
let $K$ be a connected component of $H$. Then $K$ has at most two vertices of
degree one.
\end{lemma}

\begin{proof}
Let $K$ have $n$ vertices, $m$ edges, and $d$ vertices of degree one.
Connectedness gives $m \ge n - 1$. Every vertex has degree at most two, and
the $d$ special ones have degree exactly one, so
\[
  2m \;=\; \sum_{v \in K} \deg v \;\le\; 2n - d .
\]
Combining, $2(n-1) \le 2m \le 2n - d$, whence $d \le 2$.
\end{proof}

In a properly $3$-edge-coloured cubic graph every vertex meets exactly one
edge of each colour, so a two-coloured subgraph has degree exactly two at
every internal vertex and degree one at a boundary stub. Lemma~\ref{lem:two-stubs}
therefore says a Kempe component of a colour pair contains at most two stubs
--- it is a path between them, or a cycle carrying none.

The degree hypothesis is already available in this development rather than
needing to be re-established:
\texttt{\scriptsize colorPairGraph\_degree\_eq\_two\_of\_cubic\_tait}
gives degree exactly two at any cubic vertex of a Tait colouring, and
\texttt{\scriptsize cubicInterior\_colorPairGraph\_degree\_eq\_two}
specialises it to the framed model's interior. Wiring those to
Lemma~\ref{lem:two-stubs} requires only a bound at the boundary vertices and
a connectivity hypothesis at the use site; the adapter is deliberately not
written here, since no consumer yet exists and a speculative one would be a
guess at the interface its caller wants.

\begin{remark}[where this must and must not be applied]
The lemma is about the \emph{open} tangle, where stubs are genuinely of
degree one. In the closed map the same two-coloured subgraph has degree two
everywhere and its components are cycles, which may pass through arbitrarily
many former stub positions. Counting components in the open tangle and
counting them in the closed map are therefore different measurements, and a
bound established for one says nothing about the other. The distinction is
easy to lose and expensive to lose: an argument that silently moves between
the two proves nothing.
\end{remark}

\begin{remark}[what was avoided]
The usual proof of Corollary~\ref{cor:two-sided} runs: a non-bridge lies on
a cycle; the cycle is a Jordan curve; the two sides of the edge are
separated by it; hence they lie in different faces. That argument is
correct and is what a reader should picture. It is also, in a formal
development, expensive: it requires a combinatorial Jordan curve theorem
for embedded graphs, which no current library provides and which is a
substantial development in its own right. The route above replaces it with
Lemma~\ref{lem:split-merge} and one induction. The picture is the
motivation; the counting is the proof.
\end{remark}

\begin{remark}[on the order of work]
It is worth recording why the order mattered. The target,
Corollary~\ref{cor:two-sided}, is a statement about faces, and the natural
first instinct is to attack it directly: take an edge with one face on both
sides, walk around that face, and try to extract a separation. That is the
Jordan route, and in a formal setting it does not terminate in reasonable
time. Proving the \emph{more general} Theorem~\ref{thm:euler-bound} first
--- a statement about arbitrary permutations and arbitrary transpositions,
with no map, no sphere, and no edges in sight --- turns the target into
three lines of arithmetic. The generalisation is easier than the special
case precisely because it removes the structure that made the special case
tempting to visualise.
\end{remark}
